\section{Introduction}\label{sec:introduction}
A prepared experiment supplies conditional laws of future observations, not an a priori information manifold. Once the legal future tests are fixed, two histories carry the same predictive state precisely when no such test can distinguish them. Finite resources impose a second, different equivalence: an executable observer must select its next experiment, update its retained information and produce its eventual response through the same finite state. A cover of the predictive space need not be realizable by that observer. Nor need a stationary quantizer be calibrated by the law that its own acquisition decisions produce.

The problem considered here is to derive finite-resource decision laws from this actual acquired geometry. We work with a nontrivial but explicit class: known, time-homogeneous positive causal instruments having a finite separating affine family of executable tests. The hypotheses are given before a policy, codebook or invariant law is selected. A density bound measures actual acquired mass; a common grid provides a global legal cover; positive-cone comparison and robust access to a common report component connect that geometry to a realizable recursion. Finite-dimensional closure, whole-base interiority and accessibility are substantive assumptions, not consequences of positivity alone.

Our main result, Theorem~\ref{thm:main}, constructs a single globally calibrated homogeneous observer and proves, for a protected randomly timed predictive audit,
\begin{equation}\label{eq:intro-law}
 cM^{-2/d}\le\calR_{\rm hom}(n,M)-J_*(n)
       \le C\bigl(M^{-2/d}+n^{-1}\bigr).
\end{equation}
Here $J_*(n)$ optimizes over every legal full-history adaptive acquisition policy in the same experiment. The observer uses at most $M$ total retained labels, including its acquisition control, and the same table works for every horizon $n$ and every admissible preparation. The actual support dimension $d$ is derived from executable tests and report pushforwards. It is not the ambient parameter count. The orders match for $n\gtrsim M^{2/d}$ and in the average limit. The theorem does not claim a matching $n^{-1}$ lower term or a complete short-horizon classification.

The mechanism is a lift from acquired mass to joint physical--observer mixing. Interior physical states dominate a uniform positive fraction of every programmed codepoint. Applying the raw instrument to that fraction turns a common component of its acquired law into a common component of the full unnormalized output ensemble, including the next label. Robust access to a small common region extends this one-step lift to informative instruments whose globally acquired rows do not overlap in one step. Crucially, the fraction and block length do not deteriorate as the number of labels grows. This holds for every sufficiently fine local common-grid table, before calibration or optimal control is imposed.

A compact invariant-ensemble correspondence then solves calibration and acquisition control simultaneously. It includes zero-mass labels and tied optimal actions. At its fixed point, the average Bellman equation reduces excess risk to an acquired-mass Jensen loss. Nonsmooth concavity, rather than a globally bounded Hessian, yields the second-order grid rate. The new joint minorization attracts this stationary law from every preparation without inserting an artificial label erasure, resetting the physical experiment or keeping an exact posterior register. It also amplifies per-call numerical ensemble errors by a constant independent of $M$.

The checkpoint identity and the lower in \eqref{eq:intro-law} use the actual acquired law at every scored cut. Conditional projection compares each observer with its own full-history acquisition execution; optimizing afterwards gives the all-controller converse. This avoids comparing a same-controller compression bound with an unrelated unrestricted control optimum. The external audit clock is separately paid and hidden from the observer. Internally exact-deadline tasks have a different state constraint and are not silently identified with this task.

Two structurally different raw experiments verify the same theorem. A controlled finite-state Gaussian filter has a common acquired component at every retained probability. A noncommuting quantum instrument reaches a common component through a positive-probability sequence of quiet reports; arbitrary action-dependent unitaries remain allowed. Section~\ref{sec:benchmark} gives an explicit two-label Gaussian filter, its scalar fixed-point equation and a numerical-free mixing certificate. Marked affine acquired strata yield a weighted singular profile with their true branch masses. The common-task extension adds an attained calibration alternative and causal simulation deficiency, while effective compilation separates labels, program constants, workspace, physical calls and the audit clock.

The finite predictive quotient, Bayesian conditional barycenters, acquired-curvature estimate and earlier stationary fixed-point approach are inherited results \cite{GTFBayes,GTFControlled,GTFStationary}. They are restated with proofs where they carry this theorem. The new step is the budget-independent acquired-mass lift and its joint use in calibration, finite-acquisition entrance and precision.
